% !PRINTED 275
\weberpage{275}
\weq{5}{P_m, \; P_{m-1}, \; P_{m-2} \ldots P_0}

in dem Intervall von $-1$ bis $+1$ die Eigenschaften einer \gesperrt{Sturm}'schen Kette haben. Denn wenn $P_n$ und $P_{n-1}$ für irgend ein $x$ zugleich verschwänden, so müsste nach (2) auch $P_{n-2}$, folglich $P_{n-3}$ etc. bis $P_0$ verschwinden, was unmöglich ist, da $P_0=1$ ist. Demnach ist auch die Bedingung \S.\,86, 2. für jedes beliebige Intervall befriedigt.

Ist $P_{n-1}=0$, so ist nach (2) $P_n$ und $P_{n-2}$ von entgegengesetztem Vorzeichen, wie die Bedingung \S.\,86, 3. fordert.

Aus (3) folgt, dass niemals $P_n(x)$ und $P_n'(x)$ zugleich verschwinden, da sonst auch $P_{n-1}(x)$ verschwinden müsste, und endlich ergiebt sich aus (3), dass, wenn $P_n(x)$ verschwindet und $-1<x<1$ ist, $P_n'(x)$ und $P_{n-1}(x)$ dasselbe Vorzeichen haben.

Setzen wir nun in (5) für $x$ die Werthe $-1, +1$ ein, so folgt aus (4), dass für $x=-1$ in (5) lauter Zeichenwechsel stattfinden, für $x=+1$ gar kein Zeichenwechsel; und daraus folgt:

\gesperrt{Die Gleichung $P_m(x)=0$ hat $m$ reelle Wurzeln zwischen $x=-1$ und $x=+1$, und da $P_m$ vom $m^{\text{ten}}$ Grade ist, so sind dies alle Wurzeln.}

Wir können aber einen noch etwas weiter gehenden Schluss ziehen. Nehmen wir ein Intervall $\varDelta=(\alpha,\beta)$, in dem zwei und nicht mehr Wurzeln $\xi, \eta$ von $P_m=0$ enthalten sind, so gehen beim Uebergang von $\alpha$ zu $\beta$ in der Kette (5) zwei Zeichenwechsel verloren. Wir nehmen $\alpha$ so nahe an $\xi$, $\beta$ so nahe an $\eta$, dass $P_{m-1}(\alpha)$ mit $P_{m-1}(\xi)$ und $P_{m-1}(\beta)$ mit $P_{m-1}(\eta)$ von gleichem Vorzeichen ist. Wenn $P_m$ beim Durchgang durch $\xi$ vom Negativen zum Positiven geht, so geht es beim Durchgang durch $\eta$ vom Positiven zum Negativen; es ist also $P_m'(\xi)$ und folglich $P_{m-1}(\xi)$ und $P_{m-1}(\alpha)$ positiv, $P_m'(\eta)$ und folglich $P_{m-1}(\eta)$ und $P_{m-1}(\beta)$ negativ, und das Umgekehrte findet statt, wenn $P_m$ beim Durchgang durch $\xi$ vom Positiven zum Negativen geht. Man sieht also daraus, dass beim Uebergang von $\alpha$ zu $\beta$ in

\[
P_m, \quad P_{m-1}
\]

ein Zeichenwechsel verloren geht, also muss bei dem gleichen Uebergang, und folglich auch bei dem Uebergang von $\xi$ zu $\eta$ in der Kette

\[
P_{m-1}, \quad P_{m-2} \ldots P_0
\]
%CONT>>
